\chapter{实验与分析}

本章面向云边端协同数据库中的边侧单表基数估计任务，对前文提出的分层空间划分与局部统计建模方法进行实验评估。实验设计借鉴 QDSPN 工作中对实验配置、误差分布和运行开销的呈现方式，但本文的目标并不是复现 QDSPN 的动态负载与神经概率模型实验，而是验证一种轻量级统计型基数估计方法在边侧资源约束场景下的适用性。围绕这一目标，实验关注三个问题：估计误差是否低于传统直方图基线，维度升高时误差是否保持稳定，模型构建与在线推理开销是否足够小。

\section{实验设置}

\subsection{实验环境}

实验在本地开发机上完成，程序使用 Java 实现并通过 Maven 构建。由于本文方法不依赖 GPU，实验过程仅使用 CPU 与 JVM 堆内存资源，这一点与边侧设备强调低成本部署的背景相吻合。表~\ref{tab:env} 给出了本轮实验的软硬件配置。需要说明的是，代码的 Maven 编译目标为 Java 17，当前运行环境使用 OpenJDK 26；后续若在学校论文模板中统一环境描述，可将运行 JDK 固定为 17 以增强可复现性。

\begin{table}[htbp]
  \centering
  \caption{实验环境配置}
  \label{tab:env}
  \begin{tabular}{ll}
    \hline
    项目 & 配置 \\
    \hline
    设备 & MacBook Pro \\
    处理器 & Apple M4 Max，14 核 CPU（10 个性能核心，4 个能效核心） \\
    内存 & 36 GB \\
    操作系统 & macOS 26.4.1 \\
    Java 环境 & OpenJDK 26，Maven 编译目标为 Java 17 \\
    Maven & 3.9.9 \\
    GPU & 未使用 \\
    \hline
  \end{tabular}
\end{table}

\subsection{数据集与查询负载}

为了在可控条件下观察多维分布、属性相关性和查询范围对估计误差的影响，实验采用合成数据集作为主要验证对象。数据由多簇高斯分布生成，每个数据集包含 10000 条记录，维度分别为 2 和 5，记为 \texttt{synthetic\_D2} 与 \texttt{synthetic\_D5}。低维数据集用于观察算法在简单分布上的基本表现，高维数据集则更容易暴露传统独立性假设在多谓词范围查询中的误差放大问题。数据生成参数见表~\ref{tab:data-config}。

\begin{table}[htbp]
  \centering
  \caption{合成数据生成参数}
  \label{tab:data-config}
  \begin{tabular}{lll}
    \hline
    参数 & 含义 & 取值 \\
    \hline
    $N$ & 数据规模 & 10000 \\
    $D$ & 数据维度 & 2，5 \\
    $C$ & 高斯簇数量 & 5 \\
    $[lb, ub]$ & 各维属性值域 & $[1.0, 100.0]$ \\
    $\lambda_{\mu}$ & 簇中心随机参数 & 0.2 \\
    $\lambda_{\sigma}$ & 方差随机参数 & 0.2 \\
    $\lambda_{noise}$ & 噪声幅度参数 & 0.02 \\
    SNR & 信噪比 & 30 dB \\
    \hline
  \end{tabular}
\end{table}

查询负载由合取型范围查询构成。程序按照查询中心和各维半宽随机生成矩形查询，并通过精确扫描计算真实基数，作为后续误差评价的 ground truth。当前代码中保留了 \texttt{narrow}、\texttt{medium}、\texttt{wide} 三个查询组标签，每组 200 条查询。由于现有查询生成器在难以命中目标选择率区间时会自动放宽筛选条件，本章初稿暂不把这些标签解释为严格的选择率区间；选择率敏感性实验需要在修正查询生成逻辑后再填入正式结论。现有结果仍可用于主准确性对比，因为所有方法面对的是同一批查询，方法间比较关系不受该标签问题影响。

\subsection{对比方法}

本文方法记为 \texttt{ProposedCEDEstimator}，其核心流程包括网格密度划分、稀疏区域 K-means 补充聚类、HRDAG 查询路由、轴向二分递归划分以及叶子节点一维等深直方图建模。实验选择三种传统统计估计器作为对照，覆盖从无统计结构到全局直方图的轻量基线谱系，如表~\ref{tab:methods} 所示。

\begin{table}[htbp]
  \centering
  \caption{对比方法}
  \label{tab:methods}
  \begin{tabular}{lll}
    \hline
    方法 & 记号 & 说明 \\
    \hline
    本文方法 & \texttt{ProposedCEDEstimator} & 分层空间划分与叶子局部直方图建模 \\
    等宽直方图 & \texttt{GlobalEquiWidthHistogram} & 每维全局等宽直方图，采用属性独立假设 \\
    等深直方图 & \texttt{GlobalEquiDepthHistogram} & 每维全局等深直方图，采用属性独立假设 \\
    独立性估计 & \texttt{IndependenceEstimator} & 不构建统计结构，仅按属性值域比例相乘 \\
    \hline
  \end{tabular}
\end{table}

\subsection{评价指标}

准确性评价采用基数估计领域常用的 Q-Error。给定真实基数 $A$ 与估计基数 $F$，Q-Error 定义为
\begin{equation}
  Q\text{-}Error(A,F)=\frac{\max(A,F)}{\min(A,F)} .
\end{equation}
该指标的最小值为 1，数值越接近 1 表示估计越准确。为避免单一均值掩盖尾部误差，本章报告中位数、均值、90 分位数、95 分位数和最大值。运行开销方面，实验记录模型构建耗时、全部查询的推理总耗时、单查询平均推理耗时、吞吐量以及 JVM 堆内存近似增量。堆内存增量由模型构建前后的 JVM 已用堆内存差值得到，该数值会受到垃圾回收时机影响，因此更适合作为轻量性比较的近似证据，而不是严格的对象级内存剖析结果。

\section{主准确性对比}

表~\ref{tab:main-accuracy} 汇总了两个合成数据集上的总体 Q-Error 结果。可以看到，在 D=2 数据集上，本文方法的中位数 Q-Error 为 1.040，均值为 1.055，尾部误差也被控制在较低水平；全局等宽与等深直方图在中位数上尚能维持可接受表现，但 95 分位数与最大值明显更高。进入 D=5 场景后，差距被进一步放大。本文方法的中位数 Q-Error 为 1.136，而全局等宽、全局等深和独立性估计分别上升到 5.914、6.373 和 18.891。这样的变化说明，传统方法依赖的属性独立假设会随着维度提升而迅速失效，局部空间划分则能够将联合分布的复杂性限制在更小区域内，从而降低多谓词查询中的乘法误差累积。

\begin{table}[htbp]
  \centering
  \caption{总体准确性与运行指标汇总}
  \label{tab:main-accuracy}
  \begin{tabular}{llrrrrrrr}
    \hline
    数据集 & 方法 & median & mean & p90 & p95 & max & fit(ms) & infer(us/q) \\
    \hline
    D=2 & Proposed & 1.040 & 1.055 & 1.115 & 1.153 & 1.510 & 44.53 & 4.34 \\
    D=2 & EquiWidth & 1.162 & 1.293 & 1.700 & 2.131 & 5.045 & 2.50 & 1.29 \\
    D=2 & EquiDepth & 1.239 & 1.451 & 1.962 & 2.721 & 9.471 & 9.19 & 0.38 \\
    D=2 & Independence & 1.826 & 1.954 & 2.899 & 3.359 & 3.991 & 0.02 & 0.53 \\
    D=5 & Proposed & 1.136 & 1.169 & 1.348 & 1.438 & 1.996 & 47.26 & 7.66 \\
    D=5 & EquiWidth & 5.914 & 17.183 & 39.253 & 67.602 & 345.928 & 5.14 & 0.82 \\
    D=5 & EquiDepth & 6.373 & 22.184 & 47.932 & 88.197 & 424.694 & 9.57 & 0.34 \\
    D=5 & Independence & 18.891 & 27.143 & 57.588 & 79.121 & 220.881 & 0.00 & 0.19 \\
    \hline
  \end{tabular}
\end{table}

从结果分布看，本文方法的优势不只体现在平均水平上，也体现在尾部误差的收敛性上。D=5 数据集下，全局直方图方法的最大 Q-Error 达到数百量级，这意味着部分查询会出现严重低估或高估。相比之下，本文方法的最大 Q-Error 低于 2，说明分层划分后的局部直方图没有产生数量级误差。这一点对边侧查询优化尤为重要，因为迁移决策通常对最坏情形较为敏感；一旦基数估计出现极端偏差，系统可能把本应留在边侧执行的算子迁移到云端，或在边侧承担过大的扫描压力。

\section{维度变化下的鲁棒性}

维度从 2 增加到 5 后，本文方法的中位数 Q-Error 由 1.040 上升到 1.136，退化幅度较小。全局等宽直方图从 1.162 上升到 5.914，全局等深直方图从 1.239 上升到 6.373，独立性估计则从 1.826 上升到 18.891。高维场景下，范围查询命中区域通常只占联合空间的一小部分，若仍把多维选择率简单拆成各维边缘选择率的乘积，误差会被连续放大。本文方法通过宏观簇和轴向二分把相关样本聚集在较小区域内，叶子节点中的独立性近似条件相对温和，因而能在维度升高后维持较低误差。

当前版本暂不展开严格的选择率敏感性结论。原因在于查询生成逻辑存在自动放宽区间的策略，部分查询组的真实选择率与标签含义不完全一致。后续修正时，应关闭自动放宽或将无法命中目标区间的查询直接丢弃，并在实验结果中报告每个查询组的真实选择率最小值、中位数和最大值。修正后的结果可用于绘制“选择率区间--Q-Error”折线图，补充说明窄查询、中等查询和宽查询下的稳定性差异。

\section{构建与推理开销}

边侧基数估计器不仅需要准确，也需要足够轻量。表~\ref{tab:cost} 单独列出了 D=5 数据集上的构建耗时、单查询推理耗时、吞吐量和 JVM 堆内存近似增量。本文方法在 D=5 上的构建耗时为 47.26 ms，单查询平均推理耗时为 7.66 us，吞吐量约为 130514 QPS。与全局直方图相比，本文方法需要额外构建分层空间结构，因此构建阶段和推理阶段都更重一些；不过从绝对数值看，该开销仍处在毫秒级构建和微秒级推理范围内，符合边侧“低延迟、低内存”的基本要求。

\begin{table}[htbp]
  \centering
  \caption{D=5 数据集上的构建、推理与内存开销}
  \label{tab:cost}
  \begin{tabular}{lrrrr}
    \hline
    方法 & fit(ms) & infer(us/q) & QPS & JVM 堆增量 \\
    \hline
    Proposed & 47.26 & 7.66 & 130514.14 & 395.9 KB \\
    EquiWidth & 5.14 & 0.82 & 1226647.69 & 4.2 KB \\
    EquiDepth & 9.57 & 0.34 & 2964910.29 & 4.2 KB \\
    Independence & 0.00 & 0.19 & 5293993.08 & 0 B \\
    \hline
  \end{tabular}
\end{table}

这组数据揭示了一个较清晰的权衡关系。本文方法通过额外的空间划分和局部建模换取显著更低的估计误差，代价是推理耗时比全局直方图更高；然而在 D=5 数据集上，该方法的单查询耗时仍低于 10 us，模型堆内存增量约 395.9 KB。对于一次构建、多次查询的边侧数据库场景，这样的开销通常可以接受。更严格的边端验证可在后续加入 JVM 约束运行，例如使用 \texttt{-Xmx256m -Xms256m -XX:+UseSerialGC} 模拟小内存与单线程垃圾回收环境，并将受限环境下的耗时与本地开发机结果并列表达。

\section{实验小结}

本章实验表明，本文提出的分层空间划分与局部直方图建模方法在合成数据集上具有稳定的估计精度。低维场景中，本文方法相较传统直方图已有更低的尾部误差；高维场景中，传统基线受属性独立假设影响明显退化，而本文方法仍将总体中位数 Q-Error 控制在 1.136，最大 Q-Error 控制在 2 以内。运行开销方面，本文方法在 D=5 数据集上的构建耗时约 47.26 ms，单查询推理耗时约 7.66 us，JVM 堆内存近似增量约 395.9 KB。由此可见，该方法虽然比全局直方图付出更多计算与存储成本，但其绝对开销较小，能够为边侧资源受限环境下的查询优化提供可部署的基数估计基础。

当前实验仍有两个需要完善的地方。查询选择率分组需要在生成逻辑修正后重新统计，参数敏感性与消融实验也应继续补充，以便更细致地说明网格粒度、稀疏聚类数、递归划分深度和直方图桶数对结果的影响。上述补充完成后，本章可进一步形成包含准确性、鲁棒性、资源开销和模块贡献的完整实验论证。
